\section{Implementacion computacional}

La implementacion final se organiza como paquete Python. El repositorio deja de depender
de scripts monoliticos como fuente de verdad: los scripts llaman a clases del paquete.

\begin{verbatim}
src/viu_mrob_tfm/
  domain/        robot, carga, mundo, grafo
  scenarios/     escenarios reproducibles
  allocation/    contratos de asignacion y Smith-QR
  control/       campo unico, formacion, wrench, navegacion, diff-drive
  simulation/    fachada OOP: PolicyStack, motor y resultados
  simulations/   simulador historico de benchmark de almacen
  controllers/   controladores de consenso compatibles
  experiments/   runner de configs YAML y guardado de resultados
  validation/    hipotesis, estadistica, suite
\end{verbatim}

La arquitectura tiene dos capas deliberadas. La capa OOP nueva define contratos
pequenos: \texttt{BaseAllocator}, \texttt{DecisionContext}, \texttt{BaseContinuousController},
\texttt{PolicyStack} y \texttt{SimulationEngine}. La capa historica
\texttt{simulations/warehouse.py} conserva el benchmark de almacen y las politicas
comparativas usadas por la validacion. Esta convivencia evita reescribir resultados ya
auditados mientras se extraen responsabilidades.

El contrato de implementacion es: cambiar de escenario, asignador, controlador o
configuracion no debe requerir editar el motor. Por eso la configuracion de experimento
vive en YAML y los scripts solo orquestan ejecucion, agregacion y figuras.

En V7 modular, el transporte fisico queda separado en clases explicitas.
\texttt{LoadSpec} contiene masa, destino, demanda wrench y geometria rectangular.
\texttt{RigidFormation} genera slots de contacto solidarios a la carga.
\texttt{VectorialWrenchGame} resuelve el reparto acotado de esfuerzos que aproxima
$w_k^{\rm dem}$. \texttt{SingleFieldController} consume esas piezas para sumar un termino
de contacto al campo continuo sin introducir una maquina de estados nueva.

\subsection*{Mapa de implementacion por modulo}

\begin{table}[H]
\centering
\footnotesize
\begin{tabularx}{\textwidth}{@{}p{0.13\textwidth}p{0.30\textwidth}Y@{}}
\toprule
Modulo & Artefacto principal & Lectura tecnica \\
\midrule
M0 & Benchmark de almacen y politicas greedy/centralizadas. & Comparadores clasicos y limites bajo comunicacion local. \\
M1--M1.3 & Simulador \texttt{warehouse.py}. & Asignacion a cargas, pesos heterogeneos, quorum entero y tareas dinamicas. \\
M1.4 & \texttt{LoadSpec} y juego \emph{wrench}. & Capacidad efectiva y residual wrench como reemplazo de cardinalidad pura. \\
M2 & \texttt{SingleFieldController} y \texttt{RigidFormation}. & Slots de transporte rigido y seguimiento de la carga. \\
M3 & Formacion rectangular y ranking de coalicion. & Roles fisicos derivados; falta optimizador explicito de $z_{ikh}$. \\
M4 & \texttt{SingleFieldController}. & Suma de campos de tarea, formacion, obstaculo, bateria y wrench. \\
M5 & Terminos de bateria en payoff/gates H11 y G7. & Retornabilidad incluida; estaciones de carga quedan fuera del alcance actual. \\
\bottomrule
\end{tabularx}
\caption{Trazabilidad entre modulos conceptuales y artefactos del repositorio.}
\label{tab:implementacion-modular}
\end{table}
